\section{Полиномиальная интерполяция}
\label{sec:polynomial-interpolation}

\subsection{Постановка задачи и единственность}

Пусть заданы попарно различные узлы
\[
    x_0,x_1,\ldots,x_n
\]
и значения $y_i=f(x_i)$. Требуется построить многочлен $P_n$ степени не выше
$n$ такой, что
\[
    P_n(x_i)=y_i,
    \qquad i=0,\ldots,n.
\]

Такой многочлен существует и единственен. Действительно, если два многочлена
степени не выше $n$ совпадают в $n+1$ различных точках, их разность имеет
$n+1$ корень и потому тождественно равна нулю.

\subsection{Форма Лагранжа}

Интерполяционный многочлен Лагранжа:
\[
    P_n(x)=\sum_{i=0}^{n}y_i\ell_i(x),
\]
где
\[
    \ell_i(x)=
    \prod_{\substack{j=0\\j\neq i}}^{n}
    \frac{x-x_j}{x_i-x_j}.
\]
Базисные многочлены удовлетворяют
\[
    \ell_i(x_j)=\delta_{ij}.
\]
Форма Лагранжа удобна теоретически, но при добавлении нового узла требуется
перестроение всего выражения.

\subsection{Барицентрическая формула}

Для устойчивого вычисления используют барицентрическую форму. Введём веса
\[
    w_i=\frac{1}{\prod_{j\neq i}(x_i-x_j)}.
\]
Тогда при $x\neq x_i$
\[
    P_n(x)=
    \frac{\displaystyle\sum_{i=0}^{n}\frac{w_i y_i}{x-x_i}}
         {\displaystyle\sum_{i=0}^{n}\frac{w_i}{x-x_i}}.
\]
Если $x=x_i$, полагают $P_n(x)=y_i$. После предварительного вычисления весов
стоимость одного значения --- $O(n)$.

\subsection{Форма Ньютона и разделённые разности}

Многочлен Ньютона:
\[
P_n(x)=f[x_0]
+f[x_0,x_1](x-x_0)
+\cdots+
 f[x_0,\ldots,x_n]\prod_{j=0}^{n-1}(x-x_j).
\]
Разделённые разности определяются рекурсивно:
\[
    f[x_i]=f(x_i),
\]
\[
    f[x_i,\ldots,x_{i+k}]
    =
    \frac{f[x_{i+1},\ldots,x_{i+k}]-f[x_i,\ldots,x_{i+k-1}]}
         {x_{i+k}-x_i}.
\]
При добавлении нового узла требуется вычислить только один новый коэффициент.
Многочлен удобно вычислять вложенной схемой, аналогичной схеме Горнера.

Для равноотстоящих узлов разделённые разности выражаются через конечные разности,
что приводит к формулам Ньютона вперёд и назад.

\subsection{Погрешность интерполяции}

Если $f\in C^{n+1}[a,b]$, то для каждого $x\in[a,b]$ существует
$\xi=\xi(x)$ такое, что
\[
    f(x)-P_n(x)=
    \frac{f^{(n+1)}(\xi)}{(n+1)!}
    \omega_{n+1}(x),
\]
где
\[
    \omega_{n+1}(x)=\prod_{i=0}^{n}(x-x_i).
\]
Следовательно,
\[
    |f(x)-P_n(x)|
    \leq
    \frac{\max|f^{(n+1)}|}{(n+1)!}
    |\omega_{n+1}(x)|.
\]
Погрешность зависит не только от гладкости функции, но и от расположения узлов.

\subsection{Феномен Рунге и узлы Чебышёва}

При интерполяции на равномерной сетке увеличение степени не гарантирует
сходимости. Для функции
\[
    f(x)=\frac{1}{1+25x^2}
\]
на $[-1,1]$ возникают сильные колебания у концов отрезка --- феномен Рунге.

Чтобы уменьшить максимум $|\omega_{n+1}(x)|$, выбирают узлы Чебышёва:
\[
    x_k=\cos\frac{(2k+1)\pi}{2(n+1)},
    \qquad k=0,\ldots,n,
\]
с последующим линейным переносом на $[a,b]$. Они сгущаются у концов и дают почти
оптимальную устойчивость полиномиальной интерполяции.

Чувствительность интерполятора к ошибкам данных характеризуется константой
Лебега
\[
    \Lambda_n=\max_x\sum_{i=0}^{n}|\ell_i(x)|.
\]
Для равноотстоящих узлов она растёт быстро, для узлов Чебышёва --- логарифмически.

\subsection{Интерполяция Эрмита}

Если в узлах заданы значения функции и нескольких производных, строится
интерполяционный многочлен Эрмита. Например, при заданных $f(x_i)$ и $f'(x_i)$
многочлен степени не выше $2n+1$ удовлетворяет
\[
    H(x_i)=f(x_i),
    \qquad
    H'(x_i)=f'(x_i).
\]
Его можно строить формой Ньютона с кратными узлами; разделённая разность на
совпадающих узлах определяется производной:
\[
    f[\underbrace{x_i,\ldots,x_i}_{k+1}]
    =\frac{f^{(k)}(x_i)}{k!}.
\]

\subsection{Что важно сказать на экзамене}

Следует привести теорему существования и единственности, формы Лагранжа и
Ньютона и остаточный член. Принципиальный практический момент: высокая степень на
равномерных узлах может ухудшать результат; узлы Чебышёва и барицентрическая
формула существенно повышают устойчивость.

\newpage